 %% $Id: luafindfont-doc.tex 599 2022-09-02 09:48:32Z herbert $
%
\listfiles\setcounter{errorcontextlines}{100}
\documentclass[paper=a4,fontsize=11pt,DIV=14,parskip=half-,
               captions=tableabove,twoside=on]{scrartcl}
\usepackage{fontspec}
%\usepackage{libertinus-otf}
\setmainfont{AccanthisADFStdNo3}[
  UprightFont   =*-Regular,
  BoldFont      =*-Bold,
  ItalicFont    =*-Italic,
  BoldItalicFont=*-BoldItalic,
  RawFeature    = -rlig,
]
\setsansfont{GilliusADF}[
  UprightFont   =*-Regular,
  BoldFont      =*-Bold,
  ItalicFont    =*-Italic,
  BoldItalicFont=*-BoldItalic,
  RawFeature    = -rlig,
]
\setmonofont{DejaVu Sans Mono}[Scale=MatchLowercase,FakeStretch=0.9]
%\setmonofont{Anonymous Pro}[Scale=MatchLowercase,FakeStretch=0.9]

\usepackage[english]{babel}
\usepackage{scrlayer-scrpage}
\automark[section]{section}
\automark*[subsection]{}
\pagestyle{scrheadings}

{\catcode`\%=12 
 \gdef\perCent{%}
 \gdef\DoubleperCent{%%}
}

\usepackage{xcolor,fancyvrb,varioref}
\usepackage{xltabular,pdflscape}
\usepackage{hvlogos} % for Wikipedia W
\usepackage{listings}
\usepackage{hvextern}

\setkeys{hv}{moveToExampleDir,ExampleDir=Examples,showFilename=false,force,verbose}
\lstset{basicstyle=\ttfamily\small,language=[LaTeX]TeX}
\usepackage{ragged2e}

\usepackage{makeidx}\makeindex
\usepackage{hvindex}

\def\Lfile#1{\texttt{#1}\index{#1 file@\texttt{#1} file}}
\def\Lext#1{\texttt{.#1}\index{#1 file extension@\texttt{.#1} file extension}}
\def\Lcs#1{\texttt{\textbackslash#1}\index{#1@\texttt{\textbackslash#1}}}
\def\Lenv#1{\texttt{#1}\index{#1 environment@\texttt{#1} environment}}
\def\Lpack#1{\texttt{#1}\index{#1 package@\texttt{#1} package}}
\def\Lprog#1{\texttt{#1}\index{#1 program@\texttt{#1} program}}
\def\Loption#1{\texttt{#1}\index{#1@\texttt{#1} package option}}
\def\Largb#1{\emph{\{#1\}}}

%\usepackage[bibstyle=dtk]{biblatex}
%\addbibresource{\jobname.bib}

\raggedbottom
\usepackage{url}
\usepackage[colorlinks,linktocpage]{hyperref}

\makeatletter% from: doc.sty
\newcommand*\GetFileInfo[1]{%
  \def\filename{#1}%
  \def\@tempb##1 ##2 ##3\relax##4\relax{%
    \def\filedate{##1}%
    \def\fileversion{##2}%
    \def\fileinfo{##3}}%
  \edef\@tempa{\csname ver@#1\endcsname}%
  \expandafter\@tempb\@tempa\relax? ? \relax\relax}
\makeatother

%\GetFileInfo{xltabular.sty}
\usepackage{xspace}

\newcommand\FF{\texttt{luafindfont}\xspace}

\def\setVersion#1{\setVVersion#1!!}
\def\setVVersion#1=#2!!{\def\luafindfontVersion{#2}} 

\setVersion{version = 0.10}%  can be automatically changed by perl

\setkeys{hv}{cleanup={},force}

\title{Program \texttt{luafindfont} \\--\\ \normalsize \luafindfontVersion\ (\today)}
\author{\shortstack{Herbert Voß\\\small \href{mailto:hvoss@tug.org}{hvoss@tug.org}}}
\date{}
\begin{document}
\maketitle
\thispagestyle{empty}
\enlargethispage{1cm}
\tableofcontents

\vfill
Thanks to: \\
Gerhard Bachmaier;
Karl Berry;
{\setmainfont[Scale=MatchUppercase]{NotoSerifGeorgian-Medium.ttf} მამუკა ჯიბლაძე} (Mamuka Jibladze);
Rolf Niepraschk;
\clearpage

\section{Introduction}
The Lua program
\FF is a  program for searching installed fonts which are usable for \LuaLaTeX\ and \XeLaTeX.
The already existing Lua program \Lprog{luaotfload-tool} can also be used for searching, but
the search-interface is somewhat cryptic and too complex for a user, who wants only to know
if a font is present or not. And if present, where can it be found. For example:


\begin{externalDocument}[redirect,includegraphic=false,force,code,docType=sh,ext=sh,compiler=sh]{exa}
#StartVisibleMain
luaotfload-tool --find=times
#StopVisibleMain
\end{externalDocument}

\begin{externalDocument}[redirect,includegraphic=false,force,code,docType=sh,ext=sh,compiler=sh]{exa}
#StartVisibleMain
luaotfload-tool --fuzzy --find="times new roman"
#StopVisibleMain
\end{externalDocument}




But that are not \emph{all} Times fonts. Compare it with the output of \Lprog{luafindfont},
where the searchname \emph{times} must only be a part of a longer name:

\begin{externalDocument}[redirect,includegraphic=false,force,code,docType=sh,ext=sh,compiler=sh]{exa}
#StartVisibleMain
luafindfont times
#StopVisibleMain
\end{externalDocument}



\subsection{Syntax}

The \Index{syntax} is 

\verb|luafindfont [...] <font>| 

where \verb|[...]|  are optional \Index{argument}s, either in short or
long form which, of course, can be mixed:

\begin{Verbatim}[frame=lines,label=Parameter]
  parameter handling
    -h,--help
    -o,--otfinfo (default 0)
    -i,--info (default 0)
    -x,--xetex
    -n,--no-symbolic-names
    -v...          Verbosity output
    -m,--max_string (default 90)
    <(partial) fontname>  
\end{Verbatim}

For example:

\begin{verbatim}
luafindfont times 
luafindfont palatino -o 3
luafindfont -i 3 -m 50 arial
\end{verbatim}

\begin{enumerate}
\item \texttt{-m <columns>}: How many columns are used for the font with full path. This can be really long, but with e.g. \verb|-m 50| 
the length is limited to 50 characters and in the middle of the long name there will be \ldots, representing the missing characters.
\item \texttt{-o <number[character]>}: For the font with the given number from the font list the program \Lprog{otfinfo} is called
    for some generell font informations. An additional optional argument for \Lprog{otfinfo} (only one character) must follow directly the fontnumber
    (see also section~\ref{otfinfo})
\item \texttt{-i <number>}: For the font with the given number from the font list the existing other corresponding font styles are listed.
\end{enumerate}

\section{How it works}
\FF uses the data file \Lfile{luaotfloadtool-names.lua} which is created by \LuaTeX. If it is
missing it will be created by \FF itself. The name of the font should be in lowercase
characters and must not be the complete name. A ``\verb|time|'' for finding all fonts with
``time'' in its names is sufficiant for searching ``times new roman''.


\subsection{Fontname}
The fontname can only be a substring of the  complete filename. It is not case sensitive, \texttt{Myriad} and \texttt{myriad}
are the same. It is possible to combine two search strings which will act like a boolean \texttt{and}. A search with
\texttt{"Myriad \& bold"} (in quotes)  will list all files which have \texttt{myriad} \emph{and} \texttt{bold} in its name.


\subsection{Option \texttt{-o}}\label{otfinfo}
This needs an installed program \Lprog{otfinfo}, which should be the default for all \TeX\ distributions.
This program lists some useful informations of a given fontnumber. The fontnumber can be combined by an
additional character for different outputs. For example \texttt{-o 2f} lists the features of the font number 2
in the list. The order makes no difference:  \texttt{-o f2} prints  the same. If there is no output then the font has no such
informations defined!

The possible options are (i is the default):

\begin{Verbatim}[frame=lines,label=otfinfo]
  s   Report font’s supported scripts.
  f   Report font’s GSUB/GPOS features.
  z   Report font’s optical size information.
  p   Report font’s PostScript name.
  a   Report font’s family name.
  v   Report font’s version information.
  i   Report font’s names and designer/vendor info.
  g   Report font’s glyph names.
  t   Report font’s OpenType tables.
\end{Verbatim}




\begin{externalDocument}[redirect,includegraphic=false,force,code,docType=sh,ext=sh,compiler=sh]{exa}
#StartVisibleMain
luafindfont -o 2f "myriad & semibold"
#StopVisibleMain
\end{externalDocument}

if no parameter is given, then \texttt{-i} is assumed.


\subsection{Option \texttt{-i}}
For a given font number \texttt{luafindfont} will list the font styles which were found by the font search
and which correspond to the font number. This is useful if a fontname doesn't has the style in its name or
only in short form.

\begin{externalDocument}[redirect,includegraphic=false,force,code,docType=sh,ext=sh,compiler=sh]{exa}
#StartVisibleMain
luafindfont -i 4 "myriad & semibold"
#StopVisibleMain
\end{externalDocument}



For the font number 4 there also exists a Regular, Bold, and a Italic font. The font no 4 itself
is a BoldItalic version.


\subsection{Option \texttt{-x}}\label{xetex}

\XeTeX\ will find OpenType and TrueType fonts only, if they are part of the system fonts or 
found by \texttt{kpsewhich} by its filename. With the option \texttt{-x} you'll an additional
column, which shows if the font is found by \texttt{kpsewhich} (1) or not (0). If not, then
it is obvious, that you have to update the filename database of the \TeX\ distribution in
the usual way by running \texttt{texhash} or via a menu entry if using a file manager.



\begin{externalDocument}[redirect,includegraphic=false,force,code,docType=sh,ext=sh,compiler=sh]{exa}
#StartVisibleMain
luafindfont -m 55 -x PStricksdotfont
#StopVisibleMain
\end{externalDocument}


The font \texttt{myFont.otf} is found by \LuaTeX\ but not by \XeTeX.




\section{Examples}


\begin{externalDocument}[redirect,includegraphic=false,force,code,docType=sh,ext=sh,compiler=sh]{exa}
#StartVisibleMain
luafindfont segoe
#StopVisibleMain
\end{externalDocument}

\begin{externalDocument}[redirect,includegraphic=false,force,code,docType=sh,ext=sh,compiler=sh]{exa}
#StartVisibleMain
luafindfont -m 35 libertinus
#StopVisibleMain
\end{externalDocument}

\begin{externalDocument}[redirect,includegraphic=false,force,code,docType=sh,ext=sh,compiler=sh]{exa}
#StartVisibleMain
luafindfont -i 9 sego
#StopVisibleMain
\end{externalDocument}


\begin{landscape}

\begin{externalDocument}[redirect,includegraphic=false,force,code,docType=sh,ext=sh,compiler=sh]{exa}
#StartVisibleMain
luafindfont -o 6 arial
#StopVisibleMain
\end{externalDocument}


\begin{externalDocument}[redirect,includegraphic=false,force,code,docType=sh,ext=sh,compiler=sh]{exa}
#StartVisibleMain
luafindfont --no-symbolic-names minion
#StopVisibleMain
\end{externalDocument}



\section{List of all fonts}
This possible with the symbolic fontname "*" instead of a real fontname:

\begin{Verbatim}[fontsize=\small]
bash-3.2$ luafindfont "*"
\end{Verbatim}

This list maybe really long and is shortend here!

\VerbatimInput[fontsize=\small]{Examples/exa05.txt}


\end{landscape}


\printindex

\end{document}
